\section{Torus averaging and optimal transport}
\label{tor:section}

This section establishes mean-energy convergence for the actual complex
Slater determinant under every torus-invariant reference satisfying a
logarithmic moment condition. No negative-temperature large-deviation
principle is assumed. The argument identifies a finite computational
bridge: the phase-averaged logarithm is within an explicit vanishing
error of an assignment problem with relative-entropy cost.

Put $D=q+1$, $X=\mathbb P^q$, $m=Dk$, and
$N=\binom{m+q}{q}$. Let
\[
 \mathcal A_m=\{\alpha\in\mathbb N^D:|\alpha|=m\},\qquad
 \lambda_m=N^{-1}\sum_{\alpha\in\mathcal A_m}\delta_{\alpha/m}.
\]
Write $\Delta_q=\{p\in\mathbb R_+^D:\sum p_j=1\}$ and let
$\lambda$ be normalized Lebesgue measure on this simplex. The coherent
monomial basis, orthonormal for Fubini--Study probability, is
\[
 s_\alpha(z)=\sqrt{N\binom m\alpha}\,z^\alpha,
 \qquad \binom m\alpha=\frac{m!}{\alpha_0!\cdots\alpha_q!},
 \qquad \|z\|=1.
\]
Let $\Delta_k(\boldsymbol z)=\det(s_\alpha(z_i))_{i,\alpha}$ and
$E_k=-(kN)^{-1}\log\|\Delta_k\|^2$. All phases in the statements
below are independent Haar phases, one projective torus per particle.
For $p_i\in\Delta_q^\circ$, write
$z_i=(\sqrt{p_{i0}},\sqrt{p_{i1}}e^{i\theta_{i1}},\ldots,
\sqrt{p_{iq}}e^{i\theta_{iq}})$.

\subsection{An exposed coefficient controls the phase average}

\begin{lemma}[The assignment vertices]
\label{tor:vertices}
Every monomial in $\det(z_i^\alpha)_{i,\alpha\in\mathcal A_m}$
is an exposed vertex of its Newton polytope. Consequently, if
\[
 A_m(\boldsymbol p)=
 \max_{\pi:\{1,\ldots,N\}\simeq\mathcal A_m}
 \prod_{i=1}^N N\binom m{\pi(i)}p_i^{\pi(i)},
\]
then
\begin{equation}
 \log A_m(\boldsymbol p)
 \le \int\log\|\Delta_k(\boldsymbol z)\|^2\,d\boldsymbol\theta
 \le \log A_m(\boldsymbol p)+\log N!.
 \label{tor:phase-squeeze}
\end{equation}
The logarithm is integrable for every such $\boldsymbol p$.
\end{lemma}

\begin{proof}
Delete the zeroth component from each exponent and denote the resulting
integer vector by $\alpha'\in\mathbb N^q$. For a fixed assignment
$\pi$, put $v_i=\pi(i)'$. Its scalar product with an assignment $\tau$
is $\sum_i\pi(i)'\cdot\tau(i)'$. Since both lists enumerate the same
set of distinct vectors,
\[
 \sum_i\|\pi(i)'\|^2-
 \sum_i\pi(i)'\cdot\tau(i)'
 =\frac12\sum_i\|\pi(i)'-\tau(i)'\|^2.
\]
The difference is positive unless $\tau=\pi$, proving exposure.
No two determinant terms have the same exponent vector, so the
coefficient of every term is nonzero.

For completeness, the elementary vertex form of Jensen's inequality
suffices here. For fixed $\boldsymbol\theta$, substitute
$\theta_i\mapsto\theta_i+t v_i$ into the determinant. This is a
one-variable polynomial in $e^{it}$ whose highest-degree term is the
chosen assignment, uniquely. Jensen's formula bounds its logarithmic
circle mean below by the logarithm of the modulus of that leading
coefficient. Integrating $\boldsymbol\theta$ and using invariance of
Haar measure gives the lower bound for the chosen assignment, hence for
the maximum. Logarithms of nonzero trigonometric polynomials are
integrable; one can also obtain this fact by successive one-variable
factorization and Jensen. The upper bound follows from Jensen's
inequality and orthogonality of the distinct torus characters:
\[
 \int\log|\Delta_k|^2\le
 \log\int|\Delta_k|^2
 =\log\sum_\pi\prod_iN\binom m{\pi(i)}p_i^{\pi(i)}
 \le\log A_m+\log N!.
\]
\end{proof}

For $a\in\Delta_q$, $p\in\Delta_q^\circ$, use the conventions
$0\log0=0$ and
\[
 c(a,p)=D_{\rm KL}(a\Vert p)=\sum_{j=0}^q a_j\log(a_j/p_j).
\]
Define the finite assignment cost
\[
 C_m(\boldsymbol p)=\min_\pi\frac1N\sum_{i=1}^N
                         c(\pi(i)/m,p_i).
\]
This is also the Kantorovich transport cost from $\lambda_m$ to
$N^{-1}\sum_i\delta_{p_i}$: both empirical measures have $N$ equally
weighted atoms, and the doubly stochastic polytope has permutation
vertices.

\begin{proposition}[Finite assignment certificate]
\label{tor:finite-certificate}
For every $\boldsymbol p\in(\Delta_q^\circ)^N$,
\begin{equation}
 D C_m(\boldsymbol p)-\frac{\log N}{k}
                       -\frac{\log N!}{kN}
 \le\int E_k(\boldsymbol z)\,d\boldsymbol\theta
 \le D C_m(\boldsymbol p).
 \label{tor:finite-certificate-bound}
\end{equation}
In particular the error is at most $2\log N/k$.
\end{proposition}

\begin{proof}
For $a=\alpha/m$, the multinomial type probability at its own parameter
$a$ is a largest type probability. Indeed, for another type $\beta$,
\[
 \frac{\binom m\beta a^\beta}{\binom m\alpha a^\alpha}
 =\prod_j\frac{\alpha_j!}{\beta_j!}
              \alpha_j^{\beta_j-\alpha_j}\le1.
\]
If $\alpha_j=0<\beta_j$ the numerator probability is zero; otherwise
each displayed factor is at most one, by writing out the consecutive
integer factors. The powers of $m$ cancel. There are $N$ types, so
\[
 N^{-1}\le\binom m\alpha a^\alpha\le1,
 \qquad
 e^{m h(a)}\le N\binom m\alpha\le N e^{m h(a)},
 \quad h(a)=-\sum_j a_j\log a_j.
\]
Therefore
\[
 -mN C_m\le\log A_m\le -mN C_m+N\log N.
\]
Apply Lemma~\ref{tor:vertices} and divide by $-kN$, noting $m/k=D$.
\end{proof}

\fig{fig13_toric_transport}{A finite transport certificate for Slater
energy. (a) In complex dimension $q=2$ at degree $m=6$, the 28 exponent-grid
points $a=\alpha/6$ (open blue circles) are matched to the rational
squared-modulus points $p=(\alpha+\mathbf 1)/9$ (orange). The arrows show
the identity assignment, which is optimal for this deterministic example
by the monotonicity of $\log(\alpha_j+1)$. (b) The matrix of exact costs
$D_{\rm KL}(a\Vert p)$, displayed numerically, with the optimal diagonal
shown in white. The matrix order is the same in the two coordinates.
(c) The guaranteed interval width $\log N_k/k+\log(N_k!)/(kN_k)$ in
\eqref{tor:finite-certificate-bound}, for $q=1,2,3$ and
$N_k=\binom{(q+1)k+q}{q}$. This interval controls the independent
phase average of the actual logarithmic determinant energy for fixed
moduli. It is not a pointwise bound in the phases or a bound on an
attractive exponential moment.}{fig:toric-transport}

\subsection{The transport limit, including the simplex boundary}

\begin{lemma}[Convergence of the unbounded assignment cost]
\label{tor:cost-limit}
Let $\eta$ be a probability on $\Delta_q^\circ$ such that
\begin{equation}
 \int b(p)\,d\eta(p)<\infty,
 \qquad b(p)=\sum_{j=0}^q-\log p_j.
 \label{tor:log-moment}
\end{equation}
For iid $p_1,p_2,\ldots$ with law $\eta$,
\[
 C_m(p_1,\ldots,p_N)\longrightarrow
 C(\eta):=\inf_{\kappa\in\Pi(\lambda,\eta)}
                   \int c(a,p)\,d\kappa(a,p)
\]
almost surely and in $L^1$ as $m\to\infty$ (along $m=Dk$ or all
integer degrees).
\end{lemma}

\begin{proof}
The uniform lattice measures $\lambda_m$ converge weakly to $\lambda$
by Riemann sums on the simplex. The empirical measures $\eta_N$
converge weakly almost surely. For $R>0$, let
\[
 c_R(a,p)=\sum_j a_j\log a_j+
                  \sum_j a_j\min\{-\log p_j,R\}.
\]
Extend the truncated logarithm by $R$ at $p_j=0$. Then $c_R$ is
bounded and continuous on the compact product $\Delta_q^2$.
Its transport cost is continuous in both marginals. One direct proof
couples nearby marginals by couplings converging to the diagonal,
composes these with a fixed optimal coupling, and uses uniform
continuity; compactness of optimal couplings gives the reverse bound.
Thus the $c_R$ costs for $(\lambda_m,\eta_N)$ converge to the $c_R$
cost for $(\lambda,\eta)$.

The truncation error is nonnegative and independent of the coupling:
\[
 0\le c(a,p)-c_R(a,p)
 \le t_R(p):=\sum_j(-\log p_j-R)_+.
\]
Consequently the difference of the two optimal costs is between zero
and $\int t_R\,d\eta_N$, and the analogous assertion holds for
$\eta$. For every integer $R$, the strong law makes
$\int t_R\,d\eta_N\to\int t_R\,d\eta$ simultaneously almost surely.
The latter tends to zero as $R\to\infty$ by
\eqref{tor:log-moment}. This proves almost-sure convergence.
Moreover the expected truncation error equals exactly
$\int t_R\,d\eta$, for every $N$. Bounded convergence for $c_R$ costs,
followed by $R\to\infty$, proves $L^1$ convergence. This argument does
not assume a uniform positive lower bound on the sample coordinates.
\end{proof}

\subsection{Identification with anticanonical pluricomplex energy}

\begin{proposition}[Relative-entropy transport formula]
\label{tor:energy-transport}
Let $\mu$ be invariant under the standard projective torus, and let
$\eta$ be its squared-modulus pushforward. Under
\eqref{tor:log-moment}, its anticanonical pluricomplex energy is
\begin{equation}
 E(\mu)=D C(\eta).
 \label{tor:energy-transport-formula}
\end{equation}
\end{proposition}

\begin{proof}
Averaging any admissible potential over the compact torus increases its
concave Monge--Amp\`ere energy and leaves its integral against $\mu$
unchanged. The energy dual may therefore be restricted to invariant
potentials. In logarithmic coordinates
$y_j=\log(p_j/p_0)$ put
\[
 \phi_0(y)=D\log(1+e^{y_1}+\cdots+e^{y_q}),\qquad
 \phi_0^*(Da')=D\sum_{j=0}^q a_j\log a_j,
\]
where $a'=(a_1,\ldots,a_q)$ and the moment polytope is $D\Delta_q$.
For an invariant positive metric $\phi=\phi_0+u$, normalized
Monge--Amp\`ere energy is
\[
 \mathcal E_L(u)=-\int_{\Delta_q}
                   [\phi^*(Da')-\phi_0^*(Da')]\,d\lambda(a).
\]
This follows first for smooth metrics by differentiating the Legendre
transform and changing variables by the moment map, then for finite
energy by monotone approximation. The energy dual consequently equals
\[
 \int\phi_0(y)\,d\eta_y(y)+\int\phi_0^*(Da')\,d\lambda(a)
 -\inf_\phi\left\{\int\phi(y)\,d\eta_y(y)
                   +\int\phi^*(Da')\,d\lambda(a)\right\}.
\]
The infimum is the maximal coupling integral
$\sup_{\pi\in\Pi(\lambda,\eta_y)}\int D a'\cdot y\,d\pi$.
Here is an explicit approximation argument for this duality.  First
assume that $\eta_y$ is supported on a compact cube.  Compact
Kantorovich duality for the continuous payoff $D a'\cdot y$ supplies
a bounded continuous dual function $h$ on $\Delta_q$.  Its transform
\[
 \phi(y)=\sup_{a\in\Delta_q}\{D a'\cdot y-h(a)\}
\]
has slopes in $D\Delta_q$, and $\phi^*(Da')\le h(a)$.
Fenchel's inequality gives the reverse comparison with every
coupling, so replacing the dual pair by $(\phi,\phi^*)$ attains
the same infimum.  This transform differs by a bounded function from
$D\max(0,y_1,\ldots,y_q)$, hence also from $\phi_0$.
Continuity of $h$ shows that $u=\phi-\phi_0$ extends continuously
across each coordinate face: in the limit only exponents supported
on that face contribute to the supremum.  It therefore defines a
continuous toric positive metric on $\mathbb P^q$.  Unitary heat
averaging, followed by a vanishing convex mixture with the reference
metric, approximates $u$ uniformly by smooth strictly positive toric
metrics.  Uniform convergence gives convergence both of its integrals
and of the Monge--Amp\`ere energy, so these transforms are legitimate
in the energy dual.

For general $\eta_y$, clip $y$ coordinatewise to $[-R,R]^q$.
The logarithmic moment assumption implies $\int|y|\,d\eta_y<\infty$,
so the mean clipping displacement tends to zero.  After normalizing
$\phi(0)=0$, every convex potential with slopes in $D\Delta_q$ is
$D$-Lipschitz.  Its integral changes under clipping by at most
$D\int|y-y^{(R)}|\,d\eta_y$, uniformly over all these potentials.
The same bound holds for every coupling payoff $D a'\cdot y$;
couplings can be pushed forward under clipping or lifted by
disintegration.  Thus both sides of the compact duality converge,
proving it for $\eta_y$ without any boundary-support assumption
beyond \eqref{tor:log-moment}.  The resulting cost is
\[
 \phi_0(y)+\phi_0^*(Da')-Da'\cdot y
 =D\sum_j a_j\log(a_j/p_j).
\]
This proves \eqref{tor:energy-transport-formula}. The Legendre--moment-map
correspondence is the toric optimal-transport description of
\cite{BermanWittNystrom2014}; the energy dual is part of the
variational framework of \cite{BBGZ2013}.
\end{proof}

\begin{theorem}[Actual Slater mean energy for toric reference laws]
\label{tor:mean-energy}
Let $\mu$ satisfy the hypotheses of Proposition~\ref{tor:energy-transport}.
Then, for iid particles with law $\mu$,
\begin{equation}
 \lim_{k\to\infty}\mathbb E_{\mu^{\otimes N_k}}E_k=E(\mu).
 \label{tor:mean-energy-limit}
\end{equation}
Whenever the deterministic continuous-envelope lower comparison is
available for this polarization, $E_k\to E(\mu)$ also in $L^1$ under
these product laws. In particular, for Haar probability on $\mathbb P^q$,
\[
 \mathbb E_{\sigma^{\otimes N_k}}|E_k|\longrightarrow0.
\]
On $\mathbb P^2$ this remains true after conditioning on any fixed-order
projective moment tube of width $N^{-a}$, $0<a<1/2$.
\end{theorem}

\begin{proof}
Torus invariance disintegrates $\mu$ into $\eta$ and independent uniform
angular variables. Integrate \eqref{tor:finite-certificate-bound} and
apply Lemma~\ref{tor:cost-limit} and
Proposition~\ref{tor:energy-transport}. The error tends to zero since
$\log N/k\to0$. All logarithmic determinant integrals are finite, by
the upper bound and the universal lower bound $E_k\ge-\log N/k$.

For the $L^1$ assertion, empirical measures converge to $\mu$, and the
continuous-envelope lower comparison gives
$\Pr(E_k<E(\mu)-\epsilon)\to0$ for every $\epsilon>0$.
The negative deviation $(E(\mu)-E_k)_+$ is bounded by
$E(\mu)+\log N/k$, hence tends to zero in mean. The already proved
convergence of the means then forces the positive deviation to vanish
in mean too.

For Haar probability, $\eta=\lambda$ and the diagonal coupling has
cost zero. The universal lower bound alone therefore proves $L^1$
convergence in every dimension. If $\Omega_N$ is the stated tube,
$\sigma^{\otimes N}(\Omega_N)\to1$ by the finite-moment variance
estimate. The conditional expectation of the nonnegative variable
$E_k+\log N/k$ is at most its unconditioned expectation divided by
$\sigma^{\otimes N}(\Omega_N)$, and hence also tends to zero.
\end{proof}

\paragraph{From mean energy to attractive pressure.}
The finite certificate gives a product-law mean-energy statement for the
actual Slater Hamiltonian using an ordinary assignment solver, known
constants, and explicit logarithmic boundary control. The full attractive
pressure in Section~\ref{att:entropy-section} requires the additional
entropy and conditional-density argument developed there. These are
different levels of control: convergence of the mean alone does not
permit replacing the determinant by its phase-averaged logarithm inside
an attractive partition function.
